\begin{exercisechunk}
\item \textbf{Three responses: compatibility and policy identification.}
\label[exercise]{ex:l19-three-response-comparisons}
\exerciseprofile{2}{2}{2}{\exproof\quad\excalculation\quad\exinterpretation}{%
Distinguish checking a preference model from identifying a policy within it.}
Let $\cY=\{a,b,c\}$.  In the Bradley--Terry model, a score function
$r:\cY\to\R$ predicts
\[
 p_r(y,y')=\sigma(r(y)-r(y')),
 \qquad \sigma(z)=\frac{1}{1+e^{-z}}.
\]
\begin{enumerate}[(a)]
\item Set $p(a,b)=p(b,c)=p(a,c)=3/4$, with reverse probabilities $1/4$.
Show that these probabilities admit the majority ordering $a\succ b\succ c$
but have no Bradley--Terry representation. In particular, show that
$p_r(a,b)=p_r(b,c)=3/4$ and $r(b)=0$ force
\[
 (r(a),r(b),r(c))=(\log3,0,-\log3),\qquad p_r(a,c)=9/10.
\]
\item Now suppose the comparison design has support
$\{(a,b),(b,c)\}$ and $p(a,b)=p(b,c)=3/4$. Assume the preference target
satisfies Bradley--Terry; no probability for $(a,c)$ is supplied separately.
For a uniform reference policy and $\beta=1$, the Gibbs policy is
$\pi_r^\star(y)=e^{r(y)}/\sum_{z\in\cY}e^{r(z)}$.
Show that the uniquely identified Gibbs policy is
\[
 (\pi_r^\star(a),\pi_r^\star(b),\pi_r^\star(c))=(9,3,1)/13.
\]
Explain why observing $(a,c)$ is unnecessary for this identification,
and why the two observed probabilities alone cannot test the triangle
compatibility condition.
\end{enumerate}
\end{exercisechunk}
